%-------------------------
% Resume in Latex
% Tailored for Quantum Computing & Information (CQT / IIIS PhD & Research Track)
% Candidate: Thirukumaran Kathirvelan
%------------------------

\documentclass[letterpaper,10pt]{article}

\usepackage{latexsym}
\usepackage[empty]{fullpage}
\usepackage{titlesec}
\usepackage{amsmath}
\usepackage{amssymb}
\usepackage[usenames,dvipsnames]{color}
\usepackage{verbatim}
\usepackage{enumitem}
\usepackage[hidelinks]{hyperref}
\usepackage{fancyhdr}
\usepackage[english]{babel}
\usepackage{tabularx}
\input{glyphtounicode}

\pagestyle{fancy}
\fancyhf{} % clear all header and footer fields
\fancyfoot{}
\renewcommand{\headrulewidth}{0pt}
\renewcommand{\footrulewidth}{0pt}
\setlength{\footskip}{4.5pt}

% Adjust margins for a dense, elegant single page
\addtolength{\oddsidemargin}{-0.55in}
\addtolength{\evensidemargin}{-0.55in}
\addtolength{\textwidth}{1.1in}
\addtolength{\topmargin}{-.6in}
\addtolength{\textheight}{1.2in}

\urlstyle{same}

\raggedbottom
\raggedright
\setlength{\tabcolsep}{0in}

% Sections formatting
\titleformat{\section}{
  \vspace{-5pt}\scshape\raggedright\large
}{}{0em}{}[\color{black}\titlerule \vspace{-4pt}]

% Ensure that generated pdf is machine readable/ATS parsable across all TeX engines
\ifdefined\pdfgentounicode
  \pdfgentounicode=1
\fi

%-------------------------
% Custom commands
\newcommand{\resumeItem}[1]{
  \item\small{
    {#1 \vspace{-1.5pt}}
  }
}

\newcommand{\resumeSubheading}[4]{
  \vspace{-2pt}\item
    \begin{tabular*}{0.98\textwidth}[t]{l@{\extracolsep{\fill}}r}
      \textbf{#1} & #2 \\
      \textit{\small#3} & \textit{\small #4} \\
    \end{tabular*}\vspace{-6pt}
}

\newcommand{\resumeProjectHeading}[2]{
  \vspace{-2pt}\item
    \begin{tabular*}{0.98\textwidth}{l@{\extracolsep{\fill}}r}
      \small#1 & #2 \\
    \end{tabular*}\vspace{-6pt}
}

\newcommand{\resumeSubHeadingListStart}{\begin{itemize}[leftmargin=0.15in, label={}]}
\newcommand{\resumeSubHeadingListEnd}{\end{itemize}}
\newcommand{\resumeItemListStart}{\begin{itemize}[leftmargin=0.15in, label=\textbullet]}
\newcommand{\resumeItemListEnd}{\end{itemize}\vspace{-4pt}}

%-------------------------------------------
%%%%%%  RESUME STARTS HERE  %%%%%%%%%%%%%%%%%%%%%%%%%%%%

\begin{document}

%----------HEADING----------
\begin{center}
    \textbf{\Huge \scshape Thirukumaran Kathirvelan} \\ \vspace{2pt}
    \small +91 7019214322 $|$ \href{mailto:thirukumaran3057@gmail.com}{\underline{thirukumaran3057@gmail.com}} $|$ 
    \href{https://linkedin.com/in/kthiru44}{\underline{linkedin.com/in/kthiru44}} $|$
    \href{https://github.com/kthiru-44}{\underline{github.com/kthiru-44}} $|$
    Bengaluru, India \\ \vspace{1pt}
    \textit{\small Research Interests: Quantum Information \& Algorithms, Variational Quantum Circuits, Quantum Machine Learning}
\end{center}

%-----------EDUCATION-----------
\section{Education}
  \resumeSubHeadingListStart
    \resumeSubheading
      {PES University}{Bengaluru, India}
      {Bachelor of Technology in Computer Science and Engineering; \textbf{CGPA: 9.71 / 10.0}}{}
      \resumeItemListStart
        \resumeItem{\textbf{Honors \& Awards:} Prof. M.R. Doreswamy (MRD) Merit Scholarship (50\% tuition distinction awarded to top rankers).}
        \resumeItem{\textbf{Relevant Coursework:} Quantum Computation \& Information, Linear Algebra, Multivariable Calculus, Discrete Mathematics, Probability \& Random Processes, Design \& Analysis of Algorithms, Theory of Computation, Machine Learning.}
      \resumeItemListEnd
  \resumeSubHeadingListEnd

%-----------RESEARCH & WORK EXPERIENCE-----------
\section{Research \& Leadership Experience}
  \resumeSubHeadingListStart

    \resumeSubheading
      {PESU Research Foundation (PESU-RF)}{Bengaluru, India}
      {Research Intern -- Quantum \& Applied Machine Learning}{Oct. 2026 -- Present}
      \resumeItemListStart
        \resumeItem{Investigating clinical aging dynamics on Longitudinal Ageing Study in India (LASI) data to formulate the multi-domain \textbf{Composite Aging Index (CAI)} encompassing 60+ clinical tests across physical, cognitive, biomarker, and mental health domains.}
        \resumeItem{Implemented generative multivariate Gaussian Copula models (Sklar's Theorem) for multi-domain data augmentation in the \textbf{IEEE AgeTech Global Dataset Competition}, preserving intricate inter-feature covariance structures within strict clinical limits.}
        \resumeItem{Designing parameterized quantum circuits to extract entanglement-based feature importance weights, benchmarking quantum-derived weights against classical PCA and autoencoder scoring models.}
      \resumeItemListEnd

    \resumeSubheading
      {$|\mathrm{Q}\rangle$Forest (Quantum Computing Club, PES University)}{Bengaluru, India}
      {Technical Head}{Aug. 2025 -- Present}
      \resumeItemListStart
        \resumeItem{Leading the technical direction and research reading groups of PES University's official quantum computing student organization, mentoring 20+ undergraduate peers in quantum algorithms, quantum error correction, and quantum information theory.}
        \resumeItem{Organized hands-on technical tracks and hackathons utilizing \textbf{Qiskit} and \textbf{PennyLane}, covering Bell-CHSH inequalities, Grover's search, Quantum Phase Estimation (QPE), and Variational Quantum Eigensolvers (VQE).}
      \resumeItemListEnd

    \resumeSubheading
      {IEEE Standards Association (IEEE SA)}{Bengaluru, India}
      {Research Intern -- Experiential Program on Sustainability Standards (EPSS 2026)}{Jun. 2026 -- Oct. 2026}
      \resumeItemListStart
        \resumeItem{Formulated an 8-rule verification framework specializing \textbf{ISO 14021} to validate corporate ``Plastic Neutrality'' claims across FMCG supply chains in India, addressing systemic vulnerabilities in Extended Producer Responsibility (EPR) compliance.}
        \resumeItem{Audited public sustainability disclosures and recycling declarations across 9 major FMCG conglomerates (HUL, Nestl\'e, ITC, P\&G, etc.), identifying critical gaps including unverified cement kiln co-processing and double-counted EPR credits.}
        \resumeItem{Proposed a standardized JSON/CSV schema and serial-number-based retirement registry to eliminate double-counting ; conducted hybrid visits at Without\textsuperscript{\textregistered} (Ashaya Recyclers), Toyoda Gosei, and CIPET.}
      \resumeItemListEnd

  \resumeSubHeadingListEnd

%-----------SELECTED RESEARCH PROJECTS-----------
\section{Selected Projects}
  \resumeSubHeadingListStart

    \resumeProjectHeading
      {\textbf{Qure: Quantum Kernel Genomic Classifier} $|$ \emph{Qiskit, PennyLane, Python, Scikit-learn}}{SIH Research Track}
      \resumeItemListStart
        \resumeItem{Constructed a 16-qubit Quantum Support Vector Classifier (QSVC) utilizing a 2nd-order entangling \textbf{ZZFeatureMap} with phase gates and entangling $\mathrm{CNOT} \to P(2(\pi-x_i)(\pi-x_j)) \to \mathrm{CNOT}$ blocks to encode epistatic gene interactions into Hilbert space.}
        \resumeItem{Reduced 7,129 microarray gene expressions to a 16-gene biomarker panel using cross-validated mutual information selection; evaluated Fidelity Quantum Kernel estimation on IBM Quantum Aer statevector simulators under parameterized noise models.}
      \resumeItemListEnd

    \resumeProjectHeading
      {\textbf{Quantum Kolmogorov-Arnold Networks (QKAN) \& SineKAN} $|$ \emph{PyTorch, QML}}{Theoretical Study}
      \resumeItemListStart
        \resumeItem{Reviewed foundational literature on Kolmogorov-Arnold representation theory and Quantum Kolmogorov-Arnold Networks (QKAN), studying the theoretical substitution of classical B-spline activation edges with parameterized quantum circuits (VQCs).}
        \resumeItem{Explored 1D SineKAN architectures in PyTorch utilizing Fourier basis functions as a classical precursor, analyzing the mathematical mapping from periodic sine activations to quantum single-qubit phase rotation gates.}
      \resumeItemListEnd

  \resumeSubHeadingListEnd

%-----------TECHNICAL SKILLS & HONORS-----------
\section{Technical Skills \& Honors}
 \begin{itemize}[leftmargin=0.15in, label={}]
    \small{\item{
     \textbf{Quantum Computing}{: Qiskit, PennyLane, Cirq, IBM Quantum Experience, QAOA, VQC, Quantum Kernel Methods} \\ \vspace{1pt}
     \textbf{Machine Learning \& Math}{: PyTorch, TensorFlow, Scikit-learn, Autoencoders, PCA, Gaussian Copula, SHAP, SMT/Z3, Linear Algebra, Convex Optimization, Sklar's Theorem, Statistical Hypothesis Testing} \\ \vspace{1pt}
     \textbf{Programming \& Tools}{: Python, C, C++, Linux/UNIX, Bash, Git, \LaTeX, Docker, Jupyter} \\ \vspace{1pt}
     \textbf{Honors \& Fellowships}{: \textbf{IBM Qiskit Advocate} (Student Contributor \& QGSS Mentor); \textbf{MRD Merit Scholarship} (Top Academic Ranker, PESU); \textbf{IEEE SA EPSS Fellow} (Sustainability \& Technology Track 2026)}
    }}
 \end{itemize}

%-------------------------------------------
\end{document}
